\documentclass[11pt,a4paper]{article}

\usepackage[utf8]{inputenc}
\usepackage[T1]{fontenc}
\usepackage[italian]{babel}
\usepackage{amsmath}
\usepackage{amssymb}
\usepackage[margin=2.2cm]{geometry}
\usepackage{xcolor}
\usepackage{listings}
\usepackage{enumitem}
\usepackage{booktabs}
\usepackage{longtable}
\usepackage{array}
\usepackage{titlesec}
\usepackage[hidelinks]{hyperref}

\definecolor{codebg}{rgb}{0.96,0.96,0.96}
\definecolor{kw}{rgb}{0.0,0.2,0.6}
\definecolor{cmt}{rgb}{0.3,0.5,0.3}

\lstset{
  basicstyle=\ttfamily\scriptsize,
  backgroundcolor=\color{codebg},
  keywordstyle=\color{kw}\bfseries,
  commentstyle=\color{cmt}\itshape,
  frame=single,
  framesep=4pt,
  breaklines=true,
  showstringspaces=false,
  numbers=left,
  numberstyle=\tiny\color{gray},
  language=C,
  tabsize=4
}

\newcommand{\code}[1]{\texttt{#1}}
\newcommand{\func}[1]{\textbf{\texttt{#1}}}


\titleformat{\section}{\large\bfseries\color{kw}}{\thesection}{0.6em}{}
\titleformat{\subsection}{\normalsize\bfseries}{\thesubsection}{0.6em}{}
\titleformat{\subsubsection}{\normalsize\bfseries\itshape}{\thesubsubsection}{0.6em}{}

\setlist{itemsep=1pt,topsep=2pt}

\title{\textbf{PandOSsh --- Fase 3}\\[2pt]
\large Il Support Level: memoria virtuale, Pager, U-proc e shell\\
\normalsize Guida dettagliata riga per riga, funzione per funzione}
\author{Materiale di studio per l'interrogazione}
\date{A.A. 2025/2026}

\begin{document}
\maketitle
\tableofcontents
\bigskip
\hrule
\bigskip

% ============================================================
\section{Introduzione generale alla Fase 3}
% ============================================================

La \textbf{Phase 3} (Support Level, il quarto livello dello strato
complessivo del sistema) è costruita \emph{esclusivamente} usando le 10
syscall del Nucleus esposte dalla Fase 2 (\code{SYSCALL(...)}): non tocca
mai direttamente le strutture dati interne del kernel (PCB, ASL, ready
queue). Il suo compito è realizzare due concetti che il Nucleus non conosce
affatto:

\begin{enumerate}
\item la \textbf{memoria virtuale a paginazione} per i processi utente
  (\emph{U-proc}), con relativo gestore dei page fault (il \emph{Pager});
\item i \textbf{processi utente} veri e propri, che girano in
  \textbf{user-mode} con un proprio spazio di indirizzamento logico
  (\code{kuseg}), le loro syscall dedicate (\code{TERMINATE},
  \code{WRITETERMINAL}, \code{READTERMINAL}, \code{EXECUTE}), e --- nel caso
  specifico di questo progetto --- una \textbf{shell interattiva} che
  permette di lanciare altri programmi a runtime.
\end{enumerate}

Il codice è distribuito su tre moduli C più uno assembly, con
responsabilità nettamente separate:
\begin{center}
\begin{tabular}{@{}llp{7.7cm}@{}}
\toprule
\textbf{File} & \textbf{Ruolo} \\
\midrule
\code{initProc.c} & L'\emph{InstantiatorProcess} (\code{test}): inizializza
  le strutture del Support Level, avvia la shell, attende la sua fine, poi
  provoca l'\code{HALT} finale. \\
\code{vmSupport.c} & Lo \emph{Swap Pool}, il \emph{Pager} (gestore dei page
  fault) con rimpiazzo FIFO, la lettura/scrittura sul backing store
  (device flash), l'inizializzazione della Page Table di ogni U-proc. \\
\code{sysSupport.c} & Il \emph{general exception handler}, le syscall del
  Support Level (\code{TERMINATE}, \code{WRITETERMINAL}, \code{READTERMINAL},
  \code{EXECUTE}), la terminazione ordinata (\code{supTerminate}). \\
\code{testers/crtsi.S} & Il ``C runtime startup'' delle U-proc: codice
  assembly minimale eseguito prima di \code{main} di ogni programma utente. \\
\bottomrule
\end{tabular}
\end{center}

A questi si affianca il \code{uTLB\_RefillHandler} (fisicamente collocato in
\code{phase2/exceptions.c}, compilato condizionalmente con
\code{\#ifdef SUPPORT\_LEVEL}), perché resta comunque responsabilità del
Nucleus popolare il Pass Up Vector, anche se il contenuto dell'handler serve
strutture dati (la Page Table) definite dalla Fase 3.

% ============================================================
\section{Concetti di base: memoria virtuale a paginazione}
% ============================================================

Prima di entrare nei singoli file, è indispensabile avere chiari alcuni
concetti generali di teoria dei sistemi operativi che l'implementazione
concreta (Sezioni 3--5) dà per scontati.

\begin{description}[style=nextline,leftmargin=0pt]
\item[Indirizzo logico vs fisico] Una U-proc vede solo \textbf{indirizzi
  logici} (o \emph{virtuali}), nel proprio spazio \code{kuseg} a partire da
  \code{0x80000000}: non conosce e non può indirizzare direttamente la RAM
  reale. Ogni accesso a memoria viene tradotto dall'hardware in un
  \textbf{indirizzo fisico} (un frame reale di RAM) tramite la
  \textbf{Page Table}. Questo dà a ciascuna U-proc l'illusione di avere la
  macchina (e uno spazio di indirizzamento che parte sempre da
  \code{0x80000000}) tutta per sé, isolandola dalle altre.
\item[Pagina e frame] Lo spazio logico è diviso in blocchi di dimensione
  fissa (\textbf{pagine}, qui \code{PAGESIZE = 4096} byte); la RAM fisica è
  divisa negli stessi blocchi (\textbf{frame}). Tradurre un indirizzo
  significa trovare \emph{in quale frame fisico} risiede la pagina logica
  richiesta.
\item[Page Table] Tabella (una per processo) che mappa ogni pagina logica al
  frame fisico che la contiene, più alcuni bit di controllo: qui
  \code{V} (\emph{valid}, la pagina è presente in RAM) e \code{D}
  (\emph{dirty}, la pagina è scrivibile/è stata modificata). In PandOSsh
  ogni U-proc ha una Page Table \textbf{privata} di 32 entry
  (\code{sup\_privatePgTbl}, dentro la sua \code{support\_t}): 31 pagine
  per \code{.text}/\code{.data}, 1 per lo stack.
\item[TLB (Translation Lookaside Buffer)] Consultare la Page Table in
  memoria \emph{a ogni singolo accesso} sarebbe troppo lento: il TLB è una
  piccola \textbf{cache hardware} (qui 16 entry, si veda
  \code{"tlb-size": 16} nei file di configurazione) delle traduzioni usate
  di recente. Se la traduzione richiesta non è nel TLB si parla di
  \textbf{TLB miss}: può essere risolto rapidamente (\emph{TLB-Refill}, se
  la pagina è comunque presente in Page Table) o richiedere l'intervento
  del Pager (se la pagina non è affatto in memoria).
\item[ASID (Address Space Identifier)] Un piccolo identificatore
  (\code{entry\_hi}, $\in[1,8]$ in questo progetto) che ``etichetta'' ogni
  traduzione nel TLB con il processo a cui appartiene. Senza ASID, il TLB
  andrebbe \textbf{completamente invalidato} a ogni cambio di processo
  (altrimenti un processo potrebbe leggere/scrivere per errore le pagine
  fisiche di un altro tramite una traduzione stantia); con l'ASID, più
  traduzioni di processi diversi possono coesistere nel TLB, e l'hardware
  usa solo quelle che corrispondono al processo in esecuzione.
\item[Page fault] L'eccezione generata quando si accede a una pagina
  \emph{non presente} in Page Table (\code{V=0}). \`E il segnale che fa
  scattare il \textbf{Pager} (Sezione~\ref{sub:pager-fase3}), il gestore
  che dovrà trovarla e caricarla.
\item[Paginazione a domanda (\emph{demand paging})] Strategia per cui
  nessuna pagina viene caricata in RAM finché non è effettivamente
  richiesta: ogni U-proc parte con Page Table interamente invalida
  (Sezione~\ref{sub:init-swap-pgtbl}) e le sue pagine vengono popolate una a
  una, al primo page fault su ciascuna. Evita di caricare in memoria
  codice/dati mai usati.
\item[Backing store e swapping] Poiché la RAM fisica non basta a tenere in
  memoria \emph{tutte} le pagine di \emph{tutte} le U-proc
  contemporaneamente, serve un supporto di memorizzazione più lento ma più
  capiente (qui, un device \textbf{flash} per U-proc) su cui tenere le
  pagine non attualmente in RAM: è il \textbf{backing store}. Quando la RAM
  disponibile per la memoria virtuale (lo \textbf{Swap Pool},
  Sezione~\ref{sec:vmsupport-fase3}) è
  piena, occorre \textbf{sfrattare} (\emph{swap out}) una pagina già
  presente per farne spazio a quella richiesta, riscrivendola sul suo
  backing store se necessario.
\end{description}

Questi concetti sono \textbf{teoria generale dei sistemi operativi}, validi
ben oltre PandOSsh; le sezioni successive mostrano \emph{come} questo
progetto specifico li implementa (Swap Pool a 16 frame, algoritmo FIFO,
device flash come backing store, Page Table a 32 entry per processo).

% ============================================================
\section{\code{phase3/headers/support.h} --- costanti e variabili globali}
% ============================================================

\subsection{Costanti di memoria e Swap Pool}
\begin{lstlisting}
#define SWAP_POOL_START (RAMSTART + (OSFRAMES * PAGESIZE))
#define SWAP_POOL_SIZE  POOLSIZE

#define KUSEG_VPN_START   0x80000
#define KUSEG_STACK_VPN   0xBFFFF
#define UPROC_STACKPAGE   31
#define UPROC_TEXTPAGES   31

#define SWAP_FRAME_FREE  (-1)

#define UPROC_STATUS  (MSTATUS_MPIE_MASK | MSTATUS_MPP_U)
#define SUPPORT_STATUS (MSTATUS_MIE_MASK | MSTATUS_MPIE_MASK | MSTATUS_MPP_M)
\end{lstlisting}
\begin{itemize}
\item \code{SWAP\_POOL\_START}: la RAM fisica inizia a \code{RAMSTART}
  (\code{0x20000000}); i primi \code{OSFRAMES} (32) frame sono riservati al
  Nucleus (codice, dati, stack degli handler). Lo \textbf{Swap Pool} inizia
  subito dopo, a \code{0x20000000 + 32 x 4096 = 0x20020000}.
\item \code{SWAP\_POOL\_SIZE = POOLSIZE = UPROCMAX x 2 = 16}: il numero
  di frame fisici disponibili per lo Swap Pool è il doppio del numero
  massimo di U-proc (\code{UPROCMAX = 8}), una scelta standard di Pandos
  che garantisce che ogni U-proc possa avere in memoria almeno due pagine
  contemporaneamente senza causare sfratti immediati reciproci.
\item \code{KUSEG\_VPN\_START = 0x80000}: il Virtual Page Number della
  prima pagina logica di una U-proc (corrisponde all'indirizzo virtuale
  \code{0x80000000 = KUSEG}, spostato a destra di \code{VPNSHIFT=12} bit).
\item \code{KUSEG\_STACK\_VPN = 0xBFFFF}: il VPN della pagina di stack
  (indirizzo virtuale \code{0xBFFFF000}, appena sotto
  \code{USERSTACKTOP = 0xC0000000}).
\item \code{UPROC\_STACKPAGE = 31}: nella Page Table privata di 32 entry
  (indici 0--31), la pagina di stack occupa sempre l'\textbf{ultimo}
  indice; \code{UPROC\_TEXTPAGES = 31}: le prime 31 entry (indici 0--30)
  sono riservate a \code{.text}/\code{.data}.
\item \code{SWAP\_FRAME\_FREE = -1}: valore sentinella che, nel campo
  \code{sw\_asid} di una entry \code{swap\_t}, indica che quel frame dello
  Swap Pool è attualmente libero (nessuna pagina di nessuna U-proc vi
  risiede).
\item \code{UPROC\_STATUS}: lo stato iniziale del processore per una
  U-proc: \code{MSTATUS\_MPIE\_MASK} (interrupt precedenti abilitati, che
  diventeranno gli interrupt correnti al primo \code{LDST}) e
  \code{MSTATUS\_MPP\_U} (Machine Previous Privilege = user, cioè la U-proc
  girerà in \textbf{user-mode}).
\item \code{SUPPORT\_STATUS}: lo stato usato per i \emph{context} degli
  handler del Support Level (Pager, general exception handler): kernel-mode
  (\code{MPP\_M}) con interrupt e PLT abilitati --- questi handler, pur
  facendo parte ``concettualmente'' del livello utente, girano a un
  privilegio più alto per poter eseguire le \code{SYSCALL} del Nucleus.
\end{itemize}

\subsection{Numeri di syscall e mutex sui device}
\begin{lstlisting}
#define SUP_TERMINATE      2
#define SUP_WRITEPRINTER   3
#define SUP_WRITETERMINAL  4
#define SUP_READTERMINAL   5
#define SUP_EXECUTE        6

#define MAXSTRLEN 128

#define DEVCNT            (DEVINTNUM * DEVPERINT)
#define DEV_MUTEX_TOTAL   (DEVCNT + DEVPERINT)

#define FLASH_MUTEX(dev)    (((IL_FLASH - IL_DISK) * DEVPERINT) + (dev))
#define PRINTER_MUTEX(dev)  (((IL_PRINTER - IL_DISK) * DEVPERINT) + (dev))
#define TERMW_MUTEX(dev)    (((IL_TERMINAL - IL_DISK) * DEVPERINT) + (dev))
#define TERMR_MUTEX(dev)    (DEVCNT + (dev))
\end{lstlisting}
Le syscall del Support Level hanno codice \textbf{positivo} (in contrasto
con quelle del Nucleus, negative): \code{2}=\code{TERMINATE},
\code{4}=\code{WRITETERMINAL}, \code{5}=\code{READTERMINAL},
\code{6}=\code{EXECUTE} (\code{3}=\code{WRITEPRINTER} è definita ma non
utilizzata dai programmi di test forniti). L'array \code{devMutex} fornisce
un semaforo binario \emph{per ogni sotto-device} (48 in totale: 40 device
``principali'' su 5 linee $\times$ 8, più 8 canali di ricezione terminale
separati), per serializzare l'accesso quando più U-proc contendono lo
stesso device fisico (tipicamente il flash durante il paging, o il
terminale 0 condiviso da shell e programmi).

\subsection{Variabili globali del Support Level}
\begin{lstlisting}
extern int swapPoolSem;
extern swap_t swapPool[SWAP_POOL_SIZE];
extern int masterSemaphore;
extern int shellSemaphore;
extern int devMutex[DEV_MUTEX_TOTAL];
extern support_t supportPool[UPROCMAX];
\end{lstlisting}
\begin{itemize}
\item \code{swapPoolSem}: semaforo binario (mutex) che protegge l'intera
  Swap Pool table: garantisce che un solo page fault alla volta possa
  manipolare i frame fisici e il backing store.
\item \code{swapPool[16]}: la tabella che descrive il contenuto di ciascun
  frame fisico dello Swap Pool.
\item \code{masterSemaphore}: su cui si blocca l'\code{InstantiatorProcess}
  in attesa che la shell termini.
\item \code{shellSemaphore}: su cui si blocca la shell in attesa che la
  U-proc appena lanciata (via \code{EXECUTE}) termini.
\item \code{devMutex[48]}: mutex per sotto-device.
\item \code{supportPool[8]}: pool statico di support structure, una per
  ciascuna delle 8 possibili U-proc (indicizzato per \code{ASID-1}).
\end{itemize}

% ============================================================
\section{\code{phase3/initProc.c} --- InstantiatorProcess}
% ============================================================

\subsection{\func{getSupport} --- accesso al pool}
\begin{lstlisting}
support_t *getSupport(int asid) {
    return &supportPool[asid - 1];
}
\end{lstlisting}
Funzione banale ma centrale: dato un ASID $\in [1,8]$, restituisce
l'indirizzo della sua support structure nel pool statico. Il legame
\emph{uno-a-uno e statico} tra ASID e slot dell'array elimina qualunque
necessità di allocazione dinamica o di una tabella di lookup più complessa.

\subsection{\func{launchUproc} --- creazione e avvio di una U-proc}
\begin{lstlisting}
void launchUproc(int asid) {
    support_t *sup = getSupport(asid);

    sup->sup_asid = asid;
    initUprocPageTable(sup);

    sup->sup_exceptContext[PGFAULTEXCEPT].pc       = (memaddr) pager;
    sup->sup_exceptContext[PGFAULTEXCEPT].status   = SUPPORT_STATUS;
    sup->sup_exceptContext[PGFAULTEXCEPT].stackPtr =
        (memaddr) &sup->sup_stackTLB[499];

    sup->sup_exceptContext[GENERALEXCEPT].pc       = (memaddr) generalExceptionHandler;
    sup->sup_exceptContext[GENERALEXCEPT].status   = SUPPORT_STATUS;
    sup->sup_exceptContext[GENERALEXCEPT].stackPtr =
        (memaddr) &sup->sup_stackGen[499];

    state_t s;
    for (unsigned int i = 0; i < (STATE_T_SIZE_IN_BYTES / WORDLEN); i++)
        ((unsigned int *)&s)[i] = 0;

    s.pc_epc   = UPROCSTARTADDR;
    s.reg_sp   = USERSTACKTOP;
    s.status   = UPROC_STATUS;
    s.mie      = MIE_ALL;
    s.entry_hi = asid << ASIDSHIFT;

    SYSCALL(CREATEPROCESS, (int)&s, PROCESS_PRIO_LOW, (int)sup);
}
\end{lstlisting}
Questa funzione fa \emph{tutto il necessario} per portare in vita una nuova
U-proc, sia essa la shell (ASID 1, lanciata dall'InstantiatorProcess) o un
programma qualsiasi lanciato dalla shell tramite \code{EXECUTE} (ASID
2--8). Passi:

\begin{enumerate}
\item Recupera la support structure dedicata a quell'ASID e vi registra
  l'ASID stesso.
\item \code{initUprocPageTable(sup)}: inizializza le 32 entry della Page
  Table privata (vedi \code{vmSupport.c} più sotto) tutte come \emph{non
  presenti}: il primo accesso a qualunque pagina genererà un page fault, che
  il Pager gestirà caricandola dal backing store.
\item \textbf{Configura i due \code{context\_t}} usati da
  \code{passUpOrDie} nel Nucleus per il pass-up delle eccezioni: uno per i
  page fault (\code{PGFAULTEXCEPT}, indice 0), che punta a \code{pager}
  (in \code{vmSupport.c}), e uno per tutte le altre eccezioni
  (\code{GENERALEXCEPT}, indice 1), che punta a
  \code{generalExceptionHandler} (in \code{sysSupport.c}). Ciascuno dei due
  usa \textbf{uno stack separato} interno alla support structure
  (\code{sup\_stackTLB}, \code{sup\_stackGen}, 500 word ciascuno,
  puntatore inizializzato all'ultima word valida \code{[499]}, dato che lo
  stack cresce verso il basso). \textbf{Stack distinti sono fondamentali}:
  se un page fault avviene mentre si sta ancora gestendo una syscall (e
  quindi si sta usando lo stack \code{sup\_stackGen}), usare lo stesso stack
  per il Pager ne corromperebbe il contenuto.
\item \textbf{Costruisce lo stato iniziale} \code{s} della U-proc, azzerando
  prima tutta la struttura (ciclo manuale word per word, analogo a
  \code{copyState} ma da zero anziché da una sorgente) e poi impostando:
  \begin{itemize}
  \item \code{pc\_epc = UPROCSTARTADDR} (\code{0x800000B0}): il punto di
    ingresso \emph{non} è \code{main} della U-proc, ma il simbolo
    \code{\_\_start} definito da \code{crtsi.S} (vedi più avanti), collocato
    dal linker esattamente a questo indirizzo (subito dopo l'header
    \code{.aout});
  \item \code{reg\_sp = USERSTACKTOP} (\code{0xC0000000}): cima dello stack
    utente, che coincide con l'indirizzo appena sopra la pagina di stack
    mappata all'indice 31 della Page Table;
  \item \code{status = UPROC\_STATUS}: user-mode, interrupt abilitati;
  \item \code{mie = MIE\_ALL}: tutte le sorgenti di interrupt abilitate a
    livello di maschera;
  \item \code{entry\_hi = asid << ASIDSHIFT}: registra l'ASID univoco nel
    campo alto del registro \code{EntryHI}, così che ogni traduzione TLB
    generata da questa U-proc sia automaticamente taggata con il suo ASID
    (permettendo al TLB di distinguere le pagine di U-proc diverse senza
    doverlo invalidare interamente a ogni cambio di processo).
  \end{itemize}
\item \textbf{Materializza il processo tramite il Nucleus}:
  \code{SYSCALL(CREATEPROCESS, ...)} (SYS1, la \emph{unica} syscall del
  Nucleus usata per creare processi, sia kernel sia utente) con lo stato
  appena costruito, priorità bassa, e il puntatore alla support structure
  (che rende il PCB creato una vera U-proc agli occhi del Nucleus: da questo
  momento, ogni eccezione che il Nucleus non sa gestire per questo processo
  verrà passata su tramite \code{passUpOrDie} ai due context configurati
  sopra).
\end{enumerate}

\subsection{\func{test} --- l'InstantiatorProcess}
\begin{lstlisting}
void test(void) {
    initSwapStructs();

    masterSemaphore = 0;
    shellSemaphore  = 0;
    for (int i = 0; i < DEV_MUTEX_TOTAL; i++)
        devMutex[i] = 1;

    launchUproc(1);

    SYSCALL(PASSEREN, (int)&masterSemaphore, 0, 0);

    SYSCALL(TERMPROCESS, 0, 0, 0);
}
\end{lstlisting}
\code{test} è la funzione con cui il Nucleus fa partire il \emph{primo}
processo del sistema (il suo indirizzo è impostato come \code{pc\_epc} del
PCB creato in \code{initial.c::main}). Nella build \code{phase3} questa
funzione (definita qui, non nel \code{p2test.c} usato dalla Fase 2 pura) è
l'\textbf{InstantiatorProcess}, e gira interamente in \textbf{kernel-mode}
(non ha una support struct propria: è un processo puro Fase 2). Sequenza:

\begin{enumerate}
\item \code{initSwapStructs()}: inizializza lo Swap Pool e il suo semaforo
  di mutua esclusione (dettagli in \code{vmSupport.c}).
\item Inizializza i semafori di sincronizzazione del Support Level a 0
  (\code{masterSemaphore}, \code{shellSemaphore}) e tutti i 48 mutex dei
  device a 1 (semafori binari liberi).
\item \textbf{Avvia \emph{una sola} U-proc: la shell}, con \code{ASID=1}
  (\code{launchUproc(1)}). \`E questa la scelta architetturale distintiva
  di PandOSsh rispetto allo schema ``classico'' di Pandos, che normalmente
  lancerebbe direttamente tutte le U-proc di test in sequenza: qui, invece,
  l'unico processo lanciato dal kernel è una shell interattiva, e sono le
  \emph{syscall \code{EXECUTE}} della shell stessa (invocate a runtime in
  base a cosa l'utente digita) a creare, una alla volta, le altre U-proc.
\item \textbf{Si blocca su \code{masterSemaphore}} con una \code{PASSEREN}
  (P, syscall del Nucleus): l'InstantiatorProcess resta bloccato finché la
  shell (ASID 1) non esegue una V su questo stesso semaforo --- cosa che
  accade in \code{supTerminate} quando la shell termina (comando
  \code{exit}).
\item \textbf{Sbloccato}, esegue \code{TERMPROCESS} su se stesso
  (\code{a1=0}): essendo l'InstantiatorProcess l'\emph{ultimo} processo
  rimasto vivo a questo punto (tutte le U-proc sono già terminate, essendo
  la shell l'unica lanciata direttamente e le sue figlie terminate prima di
  lei), questa terminazione porta \code{processCount} a 0, e lo scheduler
  (al successivo giro) esegue \code{HALT()}, fermando la macchina virtuale.
\end{enumerate}

% ============================================================
\section{\code{phase3/vmSupport.c} --- memoria virtuale e Pager}
\label{sec:vmsupport-fase3}
% ============================================================

\subsection{Stato del modulo}
\begin{lstlisting}
int    swapPoolSem;
swap_t swapPool[SWAP_POOL_SIZE];
static int fifoNext = 0;
\end{lstlisting}
\code{fifoNext} è l'indice (privato del modulo) usato dall'algoritmo di
rimpiazzo pagine: parte da 0 e avanza ciclicamente su tutti i 16 frame dello
Swap Pool, realizzando un semplice \textbf{round-robin} (qui usato come
approssimazione di FIFO: il frame scelto è sempre il successivo nell'ordine
circolare, indipendentemente da quando è stato effettivamente occupato ---
è comunque l'algoritmo che la specifica di Pandos ammette esplicitamente
per la sua semplicità e prevedibilità).

\subsection{Funzioni di utilità (\code{static})}

\begin{lstlisting}
static inline void interruptsOff(void) {
    setSTATUS(getSTATUS() & ~MSTATUS_MIE_MASK);
}
static inline void interruptsOn(void) {
    setSTATUS(getSTATUS() | MSTATUS_MIE_MASK);
}
\end{lstlisting}
Due wrapper minimi per disabilitare/riabilitare il bit \code{MIE} dello
\code{status} corrente: usati per rendere \emph{atomiche} le operazioni
combinate di aggiornamento Page Table + TLB (vedi sotto).

\begin{lstlisting}
static int vpnToIndex(unsigned int entryHI) {
    unsigned int vpn = (entryHI >> VPNSHIFT) & 0xFFFFF;
    if (vpn == KUSEG_STACK_VPN) return UPROC_STACKPAGE;
    if (vpn >= KUSEG_VPN_START && vpn < KUSEG_VPN_START + UPROC_TEXTPAGES)
        return (int)(vpn - KUSEG_VPN_START);
    return -1;
}
\end{lstlisting}
\func{vpnToIndex}: dato il valore del registro \code{EntryHI} salvato al
momento del page fault, estrae il VPN (bit alti, shiftati di
\code{VPNSHIFT=12}) e lo traduce nell'indice \code{[0..31]} della Page
Table privata: caso speciale per la pagina di stack (VPN
\code{0xBFFFF} $\to$ indice 31), altrimenti pagine \code{.text}/\code{.data}
(VPN in \code{[0x80000, 0x8001E]} $\to$ indici 0--30). Se il VPN richiesto
non rientra in nessuno dei due intervalli, l'indirizzo è \textbf{fuori dallo
spazio logico} della U-proc: si ritorna \code{-1}, che il chiamante
(\code{pager}) interpreta come errore fatale (program trap $\to$
terminazione).

\begin{lstlisting}
static inline memaddr frameAddr(int i) {
    return (memaddr)(SWAP_POOL_START + i * PAGESIZE);
}
\end{lstlisting}
Calcola l'indirizzo fisico del frame \code{i}-esimo dello Swap Pool (banale
aritmetica: base più \code{i} pagine).

\begin{lstlisting}
static void markPageNotValid(pteEntry_t *pte) {
    interruptsOff();
    pte->pte_entryLO &= ~VALIDON;
    TLBCLR();
    interruptsOn();
}
\end{lstlisting}
\func{markPageNotValid}: invalida un'entry di Page Table (azzera il bit
\code{V}, \code{VALIDON}) e cancella tutto il TLB (\code{TLBCLR}) in modo
\textbf{atomico} (interrupt disabilitati durante l'operazione). Usata
quando si sfratta una pagina vittima dallo Swap Pool: da questo momento
qualunque accesso a quella pagina, da parte della sua U-proc, genererà un
nuovo page fault.

\begin{lstlisting}
static void markPagePresent(pteEntry_t *pte, memaddr phys) {
    interruptsOff();
    pte->pte_entryLO = phys | DIRTYON | VALIDON;
    TLBCLR();
    setENTRYHI(pte->pte_entryHI);
    setENTRYLO(pte->pte_entryLO);
    TLBWR();
    interruptsOn();
}
\end{lstlisting}
\func{markPagePresent}: rende presente un'entry di Page Table, impostando
il frame fisico (\code{phys}) e i bit \code{D} (dirty, scrivibile) e
\code{V} (valida). \textbf{Scelta implementativa degna di nota}: oltre a
cancellare il TLB con \code{TLBCLR} (che elimina eventuali entry stantie
relative al vecchio contenuto del frame), la nuova traduzione viene scritta
\textbf{direttamente} nel TLB (\code{setENTRYHI}/\code{setENTRYLO}/\code{
TLBWR}), invece di limitarsi a invalidare e aspettare che sia un futuro
evento di TLB-Refill a ripopolarla. La specifica di Pandos ammette entrambe
le strategie (cancellare tutto il TLB, oppure individuare/riscrivere la
singola entry); questa implementazione adotta una variante della seconda,
restando quindi pienamente conforme. La motivazione pratica di questa
scelta (rispetto alla più semplice ``solo \code{TLBCLR}'') è discussa più
sotto, dopo aver introdotto il Pager: era stato osservato un
\emph{page-fault loop} in cui l'evento di TLB-Refill non veniva sempre
rigenerato correttamente subito dopo un fault gestito dal Pager, lasciando
la U-proc bloccata a ripetere indefinitamente lo stesso fault.

\subsection{\func{initSwapStructs} e \func{initUprocPageTable}}
\label{sub:init-swap-pgtbl}
\begin{lstlisting}
void initSwapStructs(void) {
    swapPoolSem = 1;
    fifoNext    = 0;
    for (int i = 0; i < SWAP_POOL_SIZE; i++)
        swapPool[i].sw_asid = SWAP_FRAME_FREE;
}
\end{lstlisting}
Inizializza il semaforo dello Swap Pool a 1 (mutex libero), azzera l'indice
FIFO, e marca tutti i 16 frame come liberi. Chiamata una sola volta
dall'InstantiatorProcess (\code{test}).

\begin{lstlisting}
void initUprocPageTable(support_t *sup) {
    int asid = sup->sup_asid;
    for (int i = 0; i < UPROC_TEXTPAGES; i++) {
        sup->sup_privatePgTbl[i].pte_entryHI =
            ((KUSEG_VPN_START + i) << VPNSHIFT) | (asid << ASIDSHIFT);
        sup->sup_privatePgTbl[i].pte_entryLO = DIRTYON;
    }
    sup->sup_privatePgTbl[UPROC_STACKPAGE].pte_entryHI =
        (KUSEG_STACK_VPN << VPNSHIFT) | (asid << ASIDSHIFT);
    sup->sup_privatePgTbl[UPROC_STACKPAGE].pte_entryLO = DIRTYON;
}
\end{lstlisting}
Per ciascuna delle 32 entry della Page Table privata, imposta
\code{pte\_entryHI} con il VPN corretto (progressivo da
\code{KUSEG\_VPN\_START} per le prime 31, fisso a \code{KUSEG\_STACK\_VPN}
per l'ultima) taggato con l'ASID della U-proc, e \code{pte\_entryLO} al solo
bit \code{DIRTYON} (scrivibile) \textbf{senza} \code{VALIDON}: ogni pagina
parte quindi \textbf{non presente} in memoria. Il primo accesso a
\emph{qualunque} pagina della U-proc genererà perciò un page fault gestito
dal Pager, che la caricherà a domanda dal backing store (paginazione
\emph{lazy}, tipica dei sistemi a domanda di pagina).

\subsection{\func{flashOperation} --- I/O sul backing store}
\begin{lstlisting}
int flashOperation(int asid, int blockNo, memaddr frameAddr_, int op) {
    int        devNo = asid - 1;
    dtpreg_t  *flash = (dtpreg_t *) DEV_REG_ADDR(IL_FLASH, devNo);
    int        mutex = FLASH_MUTEX(devNo);
    unsigned int status;

    SYSCALL(PASSEREN, (int)&devMutex[mutex], 0, 0);

    flash->data0 = frameAddr_;
    unsigned int command = ((unsigned int)blockNo << 8) | op;
    status = SYSCALL(DOIO, (int)&flash->command, (int)command, 0);

    SYSCALL(VERHOGEN, (int)&devMutex[mutex], 0, 0);
    return (int)status;
}
\end{lstlisting}
Il \textbf{backing store} di ciascuna U-proc è un device flash dedicato: la
U-proc con \code{ASID=i} usa il device flash numero \code{i-1}
(\code{devNo = asid - 1}), secondo la mappatura fissata anche nel file di
configurazione della macchina (\code{phase3\_config\_machine.json}).
\begin{enumerate}
\item Acquisisce il mutex dedicato a quel device flash (\code{PASSEREN}),
  per evitare che due page fault concorrenti (di U-proc diverse, o anche
  della stessa) tentino di usare lo stesso device contemporaneamente.
\item Imposta \code{flash->data0} all'indirizzo \code{frameAddr\_}: è il
  meccanismo \textbf{DMA} (\emph{Direct Memory Access}, trasferimento dati
  che il device esegue direttamente da/verso la RAM, senza passare
  attraverso il processore) con cui il device sa \emph{da/verso dove}
  trasferire i dati (il frame fisico coinvolto).
\item Compone il comando: il numero di blocco (\code{blockNo}, che
  corrisponde alla pagina logica, 0--31) va nei byte alti, l'opcode
  (\code{FLASHREAD} o \code{FLASHWRITE}) nel byte basso, ed esegue
  \code{DOIO} (la syscall del Nucleus) che blocca il chiamante finché
  l'operazione non è completata dall'interrupt.
\item Rilascia il mutex e restituisce lo status ricevuto.
\end{enumerate}
La mutua esclusione per-device è \textbf{deliberatamente separata} da
\code{swapPoolSem}: quest'ultimo protegge la \emph{struttura dati} condivisa
(la tabella Swap Pool), mentre \code{devMutex} protegge l'\emph{hardware}
fisico condiviso; tenerli distinti evita di serializzare inutilmente
operazioni su flash \emph{diversi} solo perché entrambe capitate durante un
page fault (in pratica, dato che c'è comunque un solo \code{swapPoolSem}
globale, un solo page fault alla volta è comunque in corso; la separazione
resta comunque corretta e più chiara concettualmente).

\subsection{\func{pager} --- il Pager, cuore della memoria virtuale}
\label{sub:pager-fase3}
\begin{lstlisting}
void pager(void) {
    support_t *sup = (support_t *) SYSCALL(GETSUPPORTPTR, 0, 0, 0);
    state_t   *exState = &sup->sup_exceptState[PGFAULTEXCEPT];

    unsigned int excCode = exState->cause & CAUSE_EXCCODE_MASK;

    if (excCode == EXC_MOD) {
        supTerminate(sup->sup_asid);
        return;
    }

    SYSCALL(PASSEREN, (int)&swapPoolSem, 0, 0);

    int p = vpnToIndex(exState->entry_hi);
    if (p < 0) {
        SYSCALL(VERHOGEN, (int)&swapPoolSem, 0, 0);
        supTerminate(sup->sup_asid);
        return;
    }

    int i = fifoNext;
    fifoNext = (fifoNext + 1) % SWAP_POOL_SIZE;
    memaddr fa = frameAddr(i);

    if (swapPool[i].sw_asid != SWAP_FRAME_FREE) {
        int        victimAsid = swapPool[i].sw_asid;
        int        victimPage = swapPool[i].sw_pageNo;
        pteEntry_t *victimPte  = swapPool[i].sw_pte;

        markPageNotValid(victimPte);

        int st = flashOperation(victimAsid, victimPage, fa, FLASHWRITE);
        if (st != READY) {
            SYSCALL(VERHOGEN, (int)&swapPoolSem, 0, 0);
            supTerminate(sup->sup_asid);
            return;
        }
    }

    int st = flashOperation(sup->sup_asid, p, fa, FLASHREAD);
    if (st != READY) {
        SYSCALL(VERHOGEN, (int)&swapPoolSem, 0, 0);
        supTerminate(sup->sup_asid);
        return;
    }

    swapPool[i].sw_asid   = sup->sup_asid;
    swapPool[i].sw_pageNo = p;
    swapPool[i].sw_pte    = &sup->sup_privatePgTbl[p];

    markPagePresent(&sup->sup_privatePgTbl[p], fa);

    SYSCALL(VERHOGEN, (int)&swapPoolSem, 0, 0);
    LDST(exState);
}
\end{lstlisting}

Il Pager è il gestore di eccezione a cui il Nucleus salta (tramite
\code{passUpOrDie(PGFAULTEXCEPT)} $\to$ \code{LDCXT}) ogni volta che una
U-proc genera un page fault vero (pagina non presente in Page Table). Gira
in kernel-mode sullo stack \code{sup\_stackTLB}. Analisi passo-passo:

\begin{enumerate}
\item \textbf{Recupero del contesto}: \code{GETSUPPORTPTR} (SYS8 del
  Nucleus) restituisce la support structure del processo che ha causato
  l'eccezione (il Nucleus, tramite \code{passUpOrDie}, aveva già copiato lo
  stato al momento del fault in \code{sup->sup\_exceptState[PGFAULTEXCEPT]}
  \emph{prima} di saltare qui: quello è \code{exState}).
\item \textbf{Caso \code{EXC\_MOD}} (TLB-Modification: tentativo di scrivere
  su una pagina segnata come read-only): in PandOSsh tutte le pagine utente
  sono marcate scrivibili (\code{DIRTYON}), quindi questa condizione non
  dovrebbe mai verificarsi in condizioni normali; se capitasse (bug applicativo
  o istruzione malformata) viene trattata come errore fatale $\to$
  \code{supTerminate}.
\item \textbf{Acquisizione del mutex} \code{swapPoolSem}: da qui in poi, un
  solo page fault alla volta può manipolare la Swap Pool table.
\item \textbf{Traduzione VPN$\to$indice} con \code{vpnToIndex}: se
  l'indirizzo è fuori dallo spazio logico valido, si rilascia il mutex e si
  termina la U-proc (program trap).
\item \textbf{Scelta del frame vittima (algoritmo FIFO)}: si prende
  l'indice corrente \code{fifoNext} e lo si fa avanzare ciclicamente,
  indipendentemente dal fatto che il frame scelto sia libero o occupato.
\item \textbf{Sfratto (se il frame era occupato)}: se \code{swapPool[i]}
  contiene già una pagina di qualche U-proc (\code{sw\_asid !=
  SWAP\_FRAME\_FREE}), quella pagina va \emph{salvata} prima di essere
  sovrascritta:
  \begin{itemize}
  \item si invalida immediatamente la sua PTE (e il TLB) con
    \code{markPageNotValid}, così che nessuno possa più accedervi
    ``credendola'' ancora presente mentre viene scritta sul disco;
  \item si scrive il contenuto del frame sul backing store \emph{della
    vittima} (\code{flashOperation(victimAsid, victimPage, fa,
    FLASHWRITE)}); se l'operazione fallisce (\code{status != READY}), si
    rilascia il mutex e si termina la U-proc \emph{corrente} (quella che ha
    causato il fault, non la vittima) come errore di I/O fatale.
  \end{itemize}
\item \textbf{Caricamento della pagina richiesta}: si legge dal backing
  store \emph{della U-proc corrente} il blocco \code{p} nel frame appena
  liberato (\code{flashOperation(sup->sup\_asid, p, fa, FLASHREAD)}); stesso
  trattamento in caso di errore.
\item \textbf{Aggiornamento della Swap Pool table} per il frame \code{i}:
  ora contiene la pagina \code{p} della U-proc corrente, con un puntatore
  diretto alla sua PTE (utile per un futuro, rapido sfratto).
\item \textbf{Aggiornamento atomico di Page Table + TLB} con
  \code{markPagePresent}: la PTE della U-proc corrente per la pagina
  \code{p} viene marcata presente, puntando al frame fisico appena
  caricato.
\item \textbf{Rilascio del mutex} e \textbf{ripresa della U-proc}
  (\code{LDST(exState)}): l'istruzione che aveva causato il fault verrà
  rieseguita da capo, questa volta trovando la traduzione valida.
\end{enumerate}

\subsection{TLB-Refill vs Pager: la distinzione più chiesta all'esame}
\label{sub:tlbrefill-vs-pager}
\begin{center}
\begin{tabular}{@{}lp{5.6cm}p{5.6cm}@{}}
\toprule
 & \textbf{TLB-Refill (\code{uTLB\_RefillHandler})} & \textbf{Pager (\code{pager})} \\
\midrule
Dove & \code{phase2/exceptions.c}, guardato da \code{\#ifdef SUPPORT\_LEVEL} &
  \code{phase3/vmSupport.c} \\
Quando scatta & la traduzione è \emph{già presente} in Page Table
  (\code{V=1}) ma non è nel TLB (cache) & la pagina è \emph{assente} in Page
  Table (\code{V=0}): non è mai stata caricata o è stata sfrattata \\
Costo & rapidissimo: solo una copia PTE$\to$TLB
  (\code{setENTRYHI}/\code{setENTRYLO}/\code{TLBWR}) & costoso: può
  richiedere di scrivere una pagina vittima su flash e di leggerne una
  nuova (due operazioni di I/O bloccanti) \\
Tocca il disco & mai & sì, se necessario (sfratto + lettura) \\
Registrato come & \code{tlb\_refill\_handler} nel Pass Up Vector (percorso
  hardware dedicato) & raggiunto via \code{passUpOrDie(PGFAULTEXCEPT)}
  $\to$ \code{LDCXT} al context configurato in \code{launchUproc} \\
\bottomrule
\end{tabular}
\end{center}

Il TLB è una \emph{cache} delle traduzioni contenute nella Page Table: può
contenere solo poche entry alla volta, quindi anche una pagina
\emph{presente} in memoria può generare un evento di TLB-Refill se la sua
traduzione non è (più) nella cache. Solo quando la pagina non è affatto in
memoria (bit \code{V=0}) si tratta di un vero page fault, di competenza del
Pager.

\subsection{Perché installare la traduzione direttamente nel TLB
(\code{markPagePresent})}
La specifica di Pandos ammette due strategie equivalenti per aggiornare il
TLB dopo che il Pager ha reso presente una pagina: (a) limitarsi a
\code{TLBCLR} e lasciare che sia un successivo evento di TLB-Refill a
ripopolare l'entry quando servirà; oppure (b) individuare/scrivere
direttamente la entry nel TLB. Questa implementazione ha scelto (una
variante di) la strategia (b) \emph{dopo aver diagnosticato un problema
concreto}: in alcune situazioni l'evento di TLB-Refill non veniva
rigenerato in modo affidabile subito dopo l'intervento del Pager, e la
U-proc restava intrappolata in un \textbf{page-fault loop} (l'accesso che
riprendeva dopo il fault ritrovava ancora una traduzione mancante nel TLB,
rigenerando lo stesso fault all'infinito). Scrivendo la traduzione
direttamente nel TLB dentro \code{markPagePresent}, l'accesso che riprende
subito dopo (\code{LDST(exState)}) trova già la entry pronta, eliminando il
problema.

% ============================================================
\section{\code{phase3/sysSupport.c} --- syscall del Support Level}
\label{sec:syssupport-fase3}
% ============================================================

\subsection{\func{supTerminate} --- terminazione ordinata di una U-proc}
\begin{lstlisting}
void supTerminate(int asid) {
    SYSCALL(PASSEREN, (int)&swapPoolSem, 0, 0);
    for (int i = 0; i < SWAP_POOL_SIZE; i++)
        if (swapPool[i].sw_asid == asid)
            swapPool[i].sw_asid = SWAP_FRAME_FREE;
    SYSCALL(VERHOGEN, (int)&swapPoolSem, 0, 0);

    if (asid == 1)
        SYSCALL(VERHOGEN, (int)&masterSemaphore, 0, 0);
    else
        SYSCALL(VERHOGEN, (int)&shellSemaphore, 0, 0);

    SYSCALL(TERMPROCESS, 0, 0, 0);
}
\end{lstlisting}
Questa è la funzione \emph{unica} tramite cui \textbf{qualunque} U-proc
smette di esistere (che sia una terminazione volontaria via \code{TERMINATE},
un errore fatale rilevato dal Pager, o un program trap gestito dal general
exception handler): mai una U-proc chiama direttamente
\code{SYSCALL(TERMPROCESS, ...)} senza passare da qui.

\begin{enumerate}
\item \textbf{Libera i frame dello Swap Pool} occupati da questa U-proc
  (sotto \code{swapPoolSem}): scorre tutta la tabella e marca come liberi
  (\code{SWAP\_FRAME\_FREE}) tutti i frame il cui \code{sw\_asid} coincide
  con l'ASID che sta terminando. \`E un passo essenziale: senza di esso, un
  futuro sfratto potrebbe tentare di scrivere sul backing store di una
  U-proc ormai terminata (e magari già sostituita da una nuova U-proc con
  lo stesso ASID, dato che gli ASID vengono riusati), corrompendone i dati.
\item \textbf{Sblocca l'attendente giusto} in base all'ASID: se la U-proc
  che termina è la \textbf{shell} (\code{ASID=1}), fa una V su
  \code{masterSemaphore} (sbloccando l'InstantiatorProcess, che porterà
  alla \code{HALT} finale); altrimenti (una qualunque U-proc \emph{figlia}
  lanciata dalla shell tramite \code{EXECUTE}) fa una V su
  \code{shellSemaphore} (sbloccando la shell, che tornerà al prompt). Questa
  distinzione realizza esplicitamente la \textbf{catena di attesa} a due
  livelli: figlia $\to$ shell $\to$ InstantiatorProcess.
\item \textbf{Terminazione effettiva} tramite il Nucleus:
  \code{SYSCALL(TERMPROCESS, 0, 0, 0)} (\code{a1=0}: termina se stesso). Da
  questo punto il PCB della U-proc viene distrutto dal Nucleus (Fase 2); la
  funzione non ritorna mai.
\end{enumerate}

\subsection{SYS4 --- \func{writeTerminal}}
\begin{lstlisting}
static int writeTerminal(support_t *sup, char *virtAddr, int len) {
    if ((memaddr)virtAddr < KUSEG || (memaddr)virtAddr >= USERSTACKTOP ||
        len < 0 || len > MAXSTRLEN) {
        supTerminate(sup->sup_asid);
    }

    termreg_t *term  = (termreg_t *) DEV_REG_ADDR(IL_TERMINAL, 0);
    int        mutex = TERMW_MUTEX(0);

    SYSCALL(PASSEREN, (int)&devMutex[mutex], 0, 0);

    int i;
    for (i = 0; i < len; i++) {
        unsigned int cmd =
            (((unsigned int)(unsigned char)virtAddr[i]) << 8) | TRANSMITCHAR;
        unsigned int status =
            SYSCALL(DOIO, (int)&term->transm_command, (int)cmd, 0);
        if ((status & 0xFF) != OKCHARTRANS) {
            SYSCALL(VERHOGEN, (int)&devMutex[mutex], 0, 0);
            return -(int)(status & 0xFF);
        }
    }

    SYSCALL(VERHOGEN, (int)&devMutex[mutex], 0, 0);
    return i;
}
\end{lstlisting}
Scrive \code{len} caratteri dal buffer \code{virtAddr} (indirizzo
\emph{logico}, nello spazio della U-proc chiamante) sul terminale 0
(condiviso da tutte le U-proc e dalla shell).
\begin{enumerate}
\item \textbf{Validazione dell'indirizzo}: deve cadere dentro lo spazio
  logico \code{[KUSEG, USERSTACKTOP)} della U-proc, e la lunghezza deve
  essere non negativa e non superiore a \code{MAXSTRLEN} (128). Se la
  richiesta è malformata, si applica la politica ``input non valido $\to$
  program trap'': la U-proc viene terminata immediatamente
  (\code{supTerminate}, che non ritorna).
\item \textbf{Mutua esclusione}: acquisisce il mutex dedicato alla
  \emph{scrittura} sul terminale 0 (\code{TERMW\_MUTEX(0)}) --- distinto dal
  mutex di lettura, perché TX e RX sono canali indipendenti e serializzarli
  insieme impedirebbe inutilmente scritture e letture concorrenti da parte
  di U-proc diverse.
\item \textbf{Ciclo di trasmissione}: per ciascun carattere, compone il
  comando (carattere nel byte alto, opcode \code{TRANSMITCHAR} nel byte
  basso) e lo emette con \code{DOIO} (bloccando finché il carattere non è
  stato trasmesso). Se lo status ricevuto non è \code{OKCHARTRANS} (5,
  successo), l'operazione si interrompe restituendo lo status come valore
  \emph{negativo} (convenzione: successo $\to$ numero di caratteri scritti,
  errore $\to$ codice di errore negato).
\item Rilascia il mutex e ritorna il numero di caratteri effettivamente
  scritti (o l'errore).
\end{enumerate}

\subsection{SYS5 --- \func{readTerminal}}
\begin{lstlisting}
static int readTerminal(support_t *sup, char *virtAddr) {
    if ((memaddr)virtAddr < KUSEG || (memaddr)virtAddr >= USERSTACKTOP) {
        supTerminate(sup->sup_asid);
    }

    termreg_t *term  = (termreg_t *) DEV_REG_ADDR(IL_TERMINAL, 0);
    int        mutex = TERMR_MUTEX(0);

    SYSCALL(PASSEREN, (int)&devMutex[mutex], 0, 0);

    int count = 0;
    int done  = 0;
    while (!done) {
        unsigned int status =
            SYSCALL(DOIO, (int)&term->recv_command, (int)RECEIVECHAR, 0);
        if ((status & 0xFF) != CHARRECV) {
            SYSCALL(VERHOGEN, (int)&devMutex[mutex], 0, 0);
            return -(int)(status & 0xFF);
        }
        char c = (char)((status >> 8) & 0xFF);
        virtAddr[count++] = c;
        if (c == '\n') done = 1;
    }

    SYSCALL(VERHOGEN, (int)&devMutex[mutex], 0, 0);
    return count;
}
\end{lstlisting}
Legge caratteri dal terminale 0 nel buffer \code{virtAddr}
\textbf{finché non incontra un \code{'\textbackslash n'}} (un vero
``leggi una riga''). Stessa validazione dell'indirizzo di
\code{writeTerminal} (qui senza controllo su una lunghezza massima, dato che
non è un parametro della syscall). Acquisisce il mutex \emph{di lettura}
(\code{TERMR\_MUTEX(0)}, indipendente da quello di scrittura). Ciclo: ad
ogni iterazione emette un comando \code{RECEIVECHAR} con \code{DOIO}
(bloccando finché non arriva un carattere), estrae il carattere ricevuto
dal byte alto dello status, lo scrive nel buffer e incrementa il contatore;
il ciclo termina quando il carattere letto è \code{'\textbackslash n'}. Se
in un qualunque momento lo status non è \code{CHARRECV} (5), l'operazione si
interrompe restituendo l'errore negato. Al termine restituisce il numero di
caratteri letti (incluso il \code{'\textbackslash n'} finale).

\subsection{SYS6 --- \func{doExecute}}
\begin{lstlisting}
static int doExecute(int asid) {
    if (asid < 1 || asid > UPROCMAX)
        return -1;
    launchUproc(asid);
    SYSCALL(PASSEREN, (int)&shellSemaphore, 0, 0);
    return 0;
}
\end{lstlisting}
Implementa la syscall con cui la shell lancia un nuovo programma: valida che
l'ASID richiesto sia nell'intervallo \code{[1, UPROCMAX]} (altrimenti
ritorna \code{-1} senza fare nulla --- \emph{non} termina il chiamante, a
differenza delle altre syscall: è un errore ``recuperabile'', gestito dalla
shell stessa che stampa un messaggio e richiede un nuovo comando), poi
chiama \code{launchUproc(asid)} (la stessa funzione usata
dall'InstantiatorProcess per lanciare la shell) e \textbf{si blocca} su
\code{shellSemaphore}: è esattamente questo blocco che rende la shell
\textbf{sincrona} --- un solo programma alla volta, il prompt torna
disponibile solo dopo che il programma appena lanciato ha eseguito
\code{supTerminate} (che, per un ASID $\ne 1$, fa la V liberatoria proprio
su \code{shellSemaphore}).

\subsection{\func{supSyscallHandler} --- dispatcher delle syscall}
\begin{lstlisting}
static void supSyscallHandler(support_t *sup, state_t *state) {
    int number = (int)state->reg_a0;
    int result = 0;

    switch (number) {
        case SUP_TERMINATE:
            supTerminate(sup->sup_asid);
            return;

        case SUP_WRITETERMINAL:
            result = writeTerminal(sup, (char *)state->reg_a1, (int)state->reg_a2);
            break;

        case SUP_READTERMINAL:
            result = readTerminal(sup, (char *)state->reg_a1);
            break;

        case SUP_EXECUTE:
            result = doExecute((int)state->reg_a1);
            break;

        default:
            supTerminate(sup->sup_asid);
            return;
    }

    state->reg_a0  = (unsigned int)result;
    state->pc_epc += WORDLEN;
    LDST(state);
}
\end{lstlisting}
Legge il numero di syscall da \code{a0} e smista alle quattro funzioni viste
sopra. Le syscall \code{TERMINATE} \emph{non ritornano mai} qui (chiamano
\code{supTerminate} che a sua volta non ritorna). Le altre tre restituiscono
un intero (\code{result}), che al termine dello \code{switch} viene scritto
in \code{a0} dello stato salvato; il PC viene avanzato di una word (per non
rieseguire la \code{ecall} originale) e si riprende la U-proc con
\code{LDST}. Qualunque numero di syscall non riconosciuto (fuori dai quattro
gestiti) viene trattato come errore fatale: \code{supTerminate}.

\subsection{\func{generalExceptionHandler} --- il punto di ingresso dal Nucleus}
\begin{lstlisting}
void generalExceptionHandler(void) {
    support_t *sup   = (support_t *) SYSCALL(GETSUPPORTPTR, 0, 0, 0);
    state_t   *state = &sup->sup_exceptState[GENERALEXCEPT];

    unsigned int excCode = state->cause & CAUSE_EXCCODE_MASK;

    if (excCode == EXC_ECU) {
        supSyscallHandler(sup, state);
    } else {
        supTerminate(sup->sup_asid);
    }
}
\end{lstlisting}
Questo è l'indirizzo configurato in \code{launchUproc} come
\code{sup\_exceptContext[GENERALEXCEPT].pc}: è dove il Nucleus salta
(tramite \code{passUpOrDie(GENERALEXCEPT)} $\to$ \code{LDCXT}) per
\emph{qualunque} eccezione non-page-fault generata da una U-proc: syscall
positive (\code{a0} $\ge 1$), program trap generici (istruzioni illegali,
accessi non allineati, ecc.). Recupera la propria support structure
(\code{GETSUPPORTPTR}) e lo stato salvato in
\code{sup\_exceptState[GENERALEXCEPT]}; distingue:
\begin{itemize}
\item se l'\code{ExcCode} è \code{EXC\_ECU} (8, \code{ecall} da user-mode):
  è davvero una richiesta di syscall $\to$ \code{supSyscallHandler};
\item altrimenti: è un vero \textbf{program trap} (errore di programma) $\to$
  terminazione immediata della U-proc con \code{supTerminate}.
\end{itemize}
Questa è la politica generale della Fase 3: \emph{qualunque} eccezione
imprevista o input malformato porta sempre e solo alla terminazione
ordinata della U-proc coinvolta, mai a un crash del sistema.

% ============================================================
\section{Il boot delle U-proc: \code{crtsi.S} e \code{libuser.h}}
% ============================================================

\subsection{\code{testers/crtsi.S} --- lo startup in assembly}
\begin{lstlisting}
#define TEXT_VADDR 0x80000000
#define TERMINATE_SYSCALL 2

    .text
    .align  2
    .globl  __start
    .type   __start, @function
    .extern main
__start:
    li  t0, TEXT_VADDR + (AOUT_HE_GP_VALUE << 2)
    lw  gp, 0(t0)

    addi    sp, sp, -16
    jal main
    addi    sp, sp, 16

    li  a0, TERMINATE_SYSCALL
    ecall
    /* Not reached! */
    .end    __start
    .size   __start, .-__start
\end{lstlisting}
Ogni programma utente in \code{testers/} (shell, echo, fibEleven, uname, sl,
calc) viene collegato con questo modulo assembly (insieme a
\code{liburiscv.S} della libreria di sistema): è il ``C runtime startup'',
l'equivalente per le U-proc di ciò che \code{crtso.S} è per il kernel. È
codice che gira interamente in \textbf{user-mode}, eseguito \emph{prima} di
\code{main}.
\begin{enumerate}
\item Definisce il simbolo \code{\_\_start}, che il linker (con lo script
  \code{uriscvaout.ldscript}) colloca esattamente all'indirizzo
  \code{0x800000B0} (\code{UPROCSTARTADDR}), cioè subito dopo l'header
  \code{.aout} che precede il \code{.text} vero e proprio (che parte a
  \code{0x80000000 = TEXT\_VADDR}). \`E questo, non \code{main}, il punto in
  cui \code{launchUproc} fa partire il PC della U-proc.
\item \textbf{Carica il \emph{global pointer}} (\code{gp}) leggendo il suo
  valore dall'header \code{.aout} (campo \code{AOUT\_HE\_GP\_VALUE},
  posizionato appena dopo \code{TEXT\_VADDR}): senza \code{gp} correttamente
  inizializzato, l'accesso a variabili globali/statiche del programma
  (indirizzate in modo relativo a \code{gp} nella convenzione RISC-V)
  fallirebbe.
\item \textbf{Alloca un piccolo frame} sullo stack (16 byte, per
  allineamento) e chiama \code{main} con \code{jal} (lo \code{sp} iniziale
  --- \code{0xC0000000} --- è già stato impostato dal kernel in
  \code{launchUproc}, prima che questo codice inizi a girare).
\item \textbf{Se \code{main} ritorna} (comportamento normale per la maggior
  parte dei programmi di test, tranne la shell che ha un loop infinito),
  imposta \code{a0 = TERMINATE\_SYSCALL} (2, cioè la stessa
  \code{SUP\_TERMINATE} del Support Level) ed esegue \code{ecall}: questa è
  a tutti gli effetti una \textbf{SYS2} (\code{TERMINATE}) che fa terminare
  la U-proc in modo pulito attraverso \code{supTerminate}. \textbf{Senza
  questo passo}, il ritorno da \code{main} (un semplice \code{ret})
  salterebbe a un indirizzo di ritorno indefinito/casuale sullo stack,
  causando un comportamento imprevedibile o un crash.
\end{enumerate}
La shell (\code{shell.c}) non sfrutta mai questa via di uscita: il suo
\code{main} contiene un ciclo infinito e termina \emph{solo} tramite una
chiamata esplicita a \code{u\_terminate()} (comando \code{exit}).

\subsection{\code{testers/h/libuser.h} --- la mini libreria utente}
\begin{lstlisting}
#define TERMINATE      2
#define WRITETERMINAL  4
#define READTERMINAL   5
#define EXECUTE        6

static void u_print(const char *s) {
    SYSCALL(WRITETERMINAL, (unsigned int)s, u_strlen(s), 0);
}
static int u_readline(char *buf) {
    return (int)SYSCALL(READTERMINAL, (unsigned int)buf, 0, 0);
}
static int u_execute(int asid) {
    return (int)SYSCALL(EXECUTE, (unsigned int)asid, 0, 0);
}
static void u_terminate(void) {
    SYSCALL(TERMINATE, 0, 0, 0);
}
\end{lstlisting}
\`E un header \emph{header-only} condiviso da tutti i programmi
\code{testers/}: fornisce wrapper minimi attorno alle quattro syscall del
Support Level
(\code{u\_print}, \code{u\_readline}, \code{u\_execute},
\code{u\_terminate}), più utility di conversione stringa/intero
(\code{u\_strlen}, \code{u\_itoa}, \code{u\_printint}), tutte scritte da
zero (non esiste una libc: le U-proc sono ``freestanding'', senza
runtime C standard).

% ============================================================
\section{La shell interattiva (\code{testers/shell.c})}
% ============================================================

\begin{lstlisting}
static int nameToAsid(const char *name) {
    if (u_strcmp(name, "echo")      == 0) return 2;
    if (u_strcmp(name, "fibEleven") == 0) return 3;
    if (u_strcmp(name, "uname")     == 0) return 4;
    if (u_strcmp(name, "sl")        == 0) return 5;
    if (u_strcmp(name, "calc")      == 0) return 6;
    return -1;
}

void main(void) {
    char buf[128];
    u_print("PandOSsh shell\n");
    u_print("comandi: echo fibEleven uname sl calc | exit\n");

    while (1) {
        u_print("> ");
        int n = u_readline(buf);
        if (n > 0 && buf[n - 1] == '\n') buf[n - 1] = '\0';
        else buf[n] = '\0';

        if (buf[0] == '\0') continue;

        if (u_strcmp(buf, "exit") == 0) {
            u_print("shell: arrivederci\n");
            u_terminate();
        }

        int asid = nameToAsid(buf);
        if (asid < 0) { u_print("shell: programma sconosciuto\n"); continue; }

        u_execute(asid);
    }
}
\end{lstlisting}
La shell (ASID 1, unica U-proc lanciata direttamente
dall'InstantiatorProcess) implementa un semplice interprete a riga di
comando:
\begin{enumerate}
\item Stampa un banner e il prompt \code{"> "}.
\item Legge una riga (\code{u\_readline}, bloccante fino a \code{'\textbackslash
  n'}), rimuove il newline finale sostituendolo con il terminatore di
  stringa.
\item Riga vuota $\to$ richiede subito un nuovo prompt (\code{continue}).
\item Comando \code{"exit"} $\to$ stampa un messaggio di commiato e chiama
  \code{u\_terminate()} (SYS2): non ritorna mai; scatena la catena
  \code{supTerminate(1)} $\to$ V su \code{masterSemaphore} $\to$
  l'InstantiatorProcess si sblocca e porta il sistema a \code{HALT}.
\item Altrimenti, traduce il nome del comando in un ASID tramite la
  mappatura statica \code{nameToAsid} (ogni programma è precaricato su un
  proprio device flash dedicato, secondo la corrispondenza \code{ASID-1 =
  numero del flash}: shell$\to$flash0, echo$\to$flash1,
  fibEleven$\to$flash2, uname$\to$flash3, sl$\to$flash4, calc$\to$flash5).
  Se il nome non è riconosciuto, avvisa e richiede un nuovo comando; se è
  valido, chiama \code{u\_execute(asid)} (SYS6), che \textbf{blocca la
  shell} finché il programma lanciato non termina, per poi tornare al
  prompt.
\end{enumerate}

% ============================================================
\section{I programmi U-proc lanciabili dalla shell (\code{testers/})}
% ============================================================

Oltre alla shell (che gira sempre come ASID 1), \code{testers/} contiene
cinque piccoli programmi utente, ciascuno collegato con lo stesso
\code{crtsi.S} e la stessa \code{libuser.h} viste sopra, e ciascuno
precaricato sul proprio device flash dedicato (\code{ASID-1 = numero del
flash}, coerentemente con \code{nameToAsid} in \code{shell.c} e con
\code{flashOperation} in \code{vmSupport.c}). Sono programmi
\emph{applicativi}, non toccano mai direttamente memoria virtuale o
Nucleus: usano esclusivamente le quattro syscall wrapper di
\code{libuser.h} (\code{u\_print}, \code{u\_readline}, \code{u\_execute},
\code{u\_terminate}). Il loro scopo è duplice: dimostrare che la
paginazione a domanda funziona per programmi realmente eseguiti (ogni
istruzione e ogni stringa che leggono risiede in pagine caricate a runtime
dal Pager) e coprire i casi richiesti esplicitamente dalla specifica di
Phase 3 (una calcolatrice, un echo, ecc.).

\subsection{\code{echo.c} --- ASID 2}
\begin{lstlisting}
#include "h/libuser.h"

void main(void) {
    char buf[128];

    u_print("echo> ");
    int n = u_readline(buf);
    if (n < 0) return;

    SYSCALL(WRITETERMINAL, (unsigned int)buf, n, 0);
}
\end{lstlisting}
Il più semplice dei cinque: stampa un prompt, legge una riga
(\code{u\_readline}, bloccante fino a \code{'\textbackslash n'}, SYS5), e se
la lettura non ha restituito un errore (\code{n < 0}, caso in cui la
funzione ritorna subito, terminando il programma tramite il consueto
epilogo di \code{crtsi.S}) la ristampa \emph{esattamente com'è}, incluso il
newline finale, con una chiamata \textbf{diretta} a
\code{SYSCALL(WRITETERMINAL, ...)} (SYS4) invece di passare per il wrapper
\code{u\_print} --- equivalente, ma dimostra che \code{u\_print} non è
altro che un caso particolare (con \code{u\_strlen} calcolato
internamente) della stessa syscall. Verifica quindi in un colpo solo sia
\code{READTERMINAL} sia \code{WRITETERMINAL}.

\subsection{\code{fibEleven.c} --- ASID 3}
\begin{lstlisting}
#include "h/libuser.h"

static int fib(int n) {
    int a = 0, b = 1;
    for (int i = 0; i < n; i++) {
        int t = a + b;
        a = b;
        b = t;
    }
    return a;
}

void main(void) {
    u_print("fib(11) = ");
    u_printint(fib(11));
    u_print("\n");
}
\end{lstlisting}
Programma puramente computazionale, senza input: calcola l'undicesimo
numero della sequenza di Fibonacci con un semplice ciclo iterativo (non
ricorsivo: evita di stressare inutilmente lo stack), poi lo stampa con
\code{u\_printint} (che internamente usa \code{u\_itoa} per convertirlo in
stringa decimale prima di passarlo a \code{u\_print}). Il nome ricorda che
in molte varianti storiche di Pandos un programma equivalente calcolava un
altro indice della sequenza; qui è fissato a 11. È il candidato ideale per
verificare che il \textbf{codice} (non solo i dati) di una U-proc venga
paginato correttamente: l'intero corpo di \code{fib} e \code{main} deve
essere caricato dal Pager (pagine \code{.text}) prima di poter essere
eseguito.

\subsection{\code{uname.c} --- ASID 4}
\begin{lstlisting}
#include "h/libuser.h"

void main(void) {
    u_print("PandOSsh - Support Level (Phase 3) - uRISCV\n");
}
\end{lstlisting}
Il programma più minimale possibile: stampa una singola stringa fissa,
imitando l'omonimo comando Unix \code{uname} (che riporta il nome del
sistema). Serve soprattutto come ``smoke test'' per verificare rapidamente
che il ciclo \code{EXECUTE} $\to$ \code{launchUproc} $\to$ page fault sul
primo accesso a \code{.text} $\to$ esecuzione $\to$ \code{TERMINATE} $\to$
\code{supTerminate} funzioni end-to-end senza alcuna altra variabile in
gioco (nessun input, nessun calcolo).

\subsection{\code{sl.c} --- ASID 5}
\begin{lstlisting}
#include "h/libuser.h"

void main(void) {
    u_print("      ====        ________                ___________\n");
    u_print("  _D _|  |_______/        \\__I_I_____===__|_________|\n");
    /* ... altre righe di ASCII art ... */
    u_print("PandOSsh locomotive\n");
}
\end{lstlisting}
Riferimento scherzoso al comando Unix \code{sl} (uno scherzo storico che
``punisce'' chi digita \code{sl} invece di \code{ls} mostrando un trenino
ASCII che attraversa lo schermo): qui è semplicemente una sequenza di
\code{u\_print} che stampa più righe di ASCII art. Dal punto di vista
tecnico non introduce nulla di nuovo rispetto a \code{uname.c}, ma essendo
un programma più lungo (più stringhe, quindi potenzialmente più pagine
\code{.text} coinvolte) esercita più a fondo il caso in cui una singola
esecuzione tocca diverse pagine di codice in sequenza.

\subsection{\code{calc.c} --- ASID 6}
\begin{lstlisting}
#include "h/libuser.h"

static void itoa(int v, char *out) {
    char tmp[16];
    int  i = 0, j = 0, neg = 0;
    if (v == 0) { out[0] = '0'; out[1] = '\0'; return; }
    if (v < 0)  { neg = 1; v = -v; }
    while (v > 0) { tmp[i++] = (char)('0' + (v % 10)); v /= 10; }
    if (neg) out[j++] = '-';
    while (i > 0) out[j++] = tmp[--i];
    out[j] = '\0';
}

void main(void) {
    char buf[128];
    char out[32];

    u_print("calc: inserisci <cifra><op><cifra> (es. 7*8): ");
    int n = u_readline(buf);

    if (n < 3) {
        u_print("calc: input non valido\n");
        return;
    }

    int  a  = buf[0] - '0';
    char op = buf[1];
    int  b  = buf[2] - '0';
    int  r  = 0;
    int  ok = 1;

    switch (op) {
        case '+': r = a + b; break;
        case '-': r = a - b; break;
        case '*': r = a * b; break;
        case '/': if (b == 0) ok = 0; else r = a / b; break;
        default:  ok = 0;
    }

    if (!ok) {
        u_print("calc: operazione non valida\n");
        return;
    }

    u_print("calc: risultato = ");
    itoa(r, out);
    u_print(out);
    u_print("\n");
}
\end{lstlisting}
Il programma applicativo più completo dei cinque, esplicitamente richiesto
dalla specifica della Phase 3: una calcolatrice a riga di comando che
accetta espressioni nella forma rigida \code{<cifra><op><cifra>} (una sola
cifra decimale per operando, un solo carattere di operatore --- non un vero
parser, ma sufficiente per lo scopo dimostrativo).
\begin{enumerate}
\item Ha una propria copia locale di \code{itoa} (identica nella logica a
  \code{u\_itoa} di \code{libuser.h}, ma duplicata invece di essere
  riutilizzata: \code{libuser.h} è header-only e ogni U-proc è compilata e
  linkata in modo del tutto indipendente dalle altre, quindi non c'è
  condivisione automatica di codice applicativo oltre a quanto già offerto
  da \code{libuser.h} stessa) usata per convertire il risultato in stringa
  prima di stamparlo.
\item Legge una riga con \code{u\_readline}; se contiene meno di 3
  caratteri (\code{n < 3}) l'input è per costruzione malformato (non può
  contenere \code{<cifra><op><cifra>}) e il programma stampa un errore e
  ritorna (terminazione pulita via \code{crtsi.S}, \emph{non} un program
  trap: qui l'errore è di logica applicativa, gestito dal programma
  stesso, non un accesso a memoria invalida).
\item Estrae i tre caratteri (\code{buf[0]}, \code{buf[1]}, \code{buf[2]}),
  converte le due cifre da carattere ASCII a valore numerico
  (\code{'0'} sottratto), ed esegue l'operazione richiesta con un normale
  \code{switch}; la divisione per zero è esplicitamente controllata
  (\code{ok = 0}) invece di lasciarla generare un errore hardware.
\item Se l'operatore non è riconosciuto o la divisione non è valida,
  stampa un messaggio di errore; altrimenti stampa il risultato
  convertito in stringa.
\end{enumerate}
Rispetto agli altri quattro programmi, \code{calc.c} è quello che meglio
dimostra l'uso \emph{combinato} di lettura (SYS5) e scrittura (SYS4) sul
terminale all'interno di un'unica esecuzione, oltre a un minimo di logica
di validazione dell'input --- lo stesso principio (``input malformato non
deve mai compromettere il sistema'') applicato a un livello più
grossolano dalle validazioni di \code{writeTerminal}/\code{readTerminal}
nel Support Level (Sezione~\ref{sec:syssupport-fase3}).

% ============================================================
\section{La pipeline di compilazione: da \code{.c} a device flash}
% ============================================================

Esistono \textbf{due build indipendenti} nel progetto:
\begin{enumerate}
\item Il \textbf{kernel} (Nucleus + Support Level): il target CMake
  \code{make phase3} compila tutti i moduli C del kernel (compresa la Fase
  3) e produce \code{build/MultiPandOS.core.uriscv} --- il file che
  l'emulatore carica come \code{core-file}, secondo
  \code{phase3\_config\_machine.json}.
\item I \textbf{programmi utente di test} (le U-proc): il \code{Makefile}
  in \code{testers/} compila ciascun programma in una propria
  \textbf{immagine flash}, indipendente dal core del kernel.
\end{enumerate}

\subsection{Il target \code{phase3} in \code{CMakeLists.txt}}
\begin{lstlisting}
add_executable(MultiPandOS_phase3 EXCLUDE_FROM_ALL
    ./klog.c
    phase1/pcb.c
    phase1/asl.c
    phase2/initial.c
    phase2/scheduler.c
    phase2/exceptions.c
    phase2/interrupts.c
    phase3/initProc.c
    phase3/vmSupport.c
    phase3/sysSupport.c
    ${URISCV_SRC}/crtso.S
    ${URISCV_SRC}/liburiscv.S
)
target_compile_definitions(MultiPandOS_phase3 PRIVATE SUPPORT_LEVEL)

add_custom_target(phase3
    COMMAND ${CMAKE_COMMAND} -E copy
        ${PROJECT_BINARY_DIR}/MultiPandOS_phase3
        ${PROJECT_BINARY_DIR}/MultiPandOS
    COMMAND uriscv-elf2uriscv -k ${PROJECT_BINARY_DIR}/MultiPandOS
    COMMAND ${CMAKE_MAKE_PROGRAM} -C ${CMAKE_SOURCE_DIR}/testers
    BYPRODUCTS MultiPandOS.core.uriscv MultiPandOS.stab.uriscv
    DEPENDS MultiPandOS_phase3
)
\end{lstlisting}
Rispetto ai target \code{phase1}/\code{phase2} (documento di Fase 1), qui
compaiono \emph{tutti e quattro} i moduli del Nucleus (\code{phase2/*.c},
\emph{non} \code{p2test.c}: il processo iniziale \code{test} viene preso da
\code{phase3/initProc.c}, non più dal test driver) più i tre moduli della
Fase 3. Due dettagli decisivi:
\begin{itemize}
\item \code{target\_compile\_definitions(\ldots PRIVATE SUPPORT\_LEVEL)}:
  è la macro di preprocessore che attiva il ramo \code{\#ifdef
  SUPPORT\_LEVEL} in \code{uTLB\_RefillHandler} (\code{phase2/exceptions.c}
  --- si veda la Sezione~\ref{sub:tlbrefill-vs-pager} di questo documento):
  \emph{solo} compilando
  questo target il TLB-Refill handler usa davvero la Page Table della
  U-proc corrente invece dello skeleton placeholder della Fase 2 pura.
\item il target \emph{custom} \code{phase3} include un
  \textbf{comando aggiuntivo} rispetto a \code{phase1}/\code{phase2}:
  \code{\$\{CMAKE\_MAKE\_PROGRAM\} -C \ldots/testers} invoca automaticamente
  anche il \code{Makefile} di \code{testers/}, così che un singolo
  \code{make phase3} produca sia il core del kernel sia le immagini flash
  di tutte le U-proc, tenendole allineate.
\end{itemize}

\subsection{\code{testers/Makefile} --- compilazione delle U-proc}
\begin{lstlisting}
CFLAGS = -ffreestanding -nostdlib -nostartfiles -Wall -O0 -std=gnu99 \
         -march=rv32imafd -mabi=ilp32d -I$(URISCV_INC) -Ih
LDFLAGS = -nostdlib -nostartfiles -march=rv32imafd -mabi=ilp32d \
          -T $(URISCV_SHARE)/uriscvaout.ldscript

PROGS = shell echo fibEleven uname sl calc
FLASHES = $(PROGS:%=%.uriscv)

all: $(FLASHES)

crtsi.o: crtsi.S
	$(CC) $(CFLAGS) -c $< -o $@

liburiscv.o: $(URISCV_SHARE)/liburiscv.S
	$(CC) $(CFLAGS) -c $< -o $@

%.o: %.c
	$(CC) $(CFLAGS) -c $< -o $@

%.elf: %.o crtsi.o liburiscv.o
	$(CC) $(LDFLAGS) crtsi.o $< liburiscv.o -o $@

%.uriscv: %.elf
	uriscv-elf2uriscv -a $<
	uriscv-mkdev -f $@ $<.aout.uriscv

clean:
	rm -f *.o *.elf *.aout.uriscv $(FLASHES)
\end{lstlisting}
\`E un Makefile classico basato su \textbf{regole implicite con pattern}
(\code{\%.o}, \code{\%.elf}, \code{\%.uriscv}): applicando la stessa
sequenza di regole a ciascun nome in \code{PROGS} (\code{shell echo
fibEleven uname sl calc}), produce automaticamente tutte e sei le immagini
flash senza dover scrivere una regola separata per ognuna. Per ciascun
programma \code{name.c}, la catena di trasformazioni è:
\begin{enumerate}
\item \textbf{\code{name.c} $\to$ \code{name.o}} (regola pattern generica
  \code{\%.o: \%.c}): compilazione con gli stessi flag ``freestanding'' del
  kernel (\code{-ffreestanding -nostdlib -nostartfiles}, nessuna libc), più
  \code{-Ih} per trovare \code{testers/h/libuser.h}.
\item \textbf{Link}: \code{name.o + crtsi.o + liburiscv.o} $\to$
  \code{name.elf} (regola \code{\%.elf: \%.o crtsi.o liburiscv.o}), con lo
  script linker \code{uriscvaout.ldscript} (diverso da
  \code{uriscvcore.ldscript} usato per il kernel: produce un formato aout
  eseguibile a partire da \code{0x80000000}, non un'immagine \code{.core}).
  \code{crtsi.o} e \code{liburiscv.o} sono \textbf{oggetti condivisi} da
  tutti i programmi (compilati una sola volta e riusati nel link di
  ciascuno): è qui che entra in gioco \code{crtsi.S}, il cui simbolo
  \code{\_\_start} finisce esattamente a \code{0x800000B0}
  (\code{UPROCSTARTADDR}).
\item \textbf{Conversione di formato}: \code{uriscv-elf2uriscv -a}
  traduce l'ELF nel formato \textbf{\code{.aout}} (load image) atteso dalla
  macchina virtuale $\to$ \code{name.elf.aout.uriscv}.
\item \textbf{Creazione dell'immagine flash}: \code{uriscv-mkdev -f
  name.uriscv name.elf.aout.uriscv} precarica il \code{.aout} su un device
  flash virtuale $\to$ \code{name.uriscv} (regola \code{\%.uriscv:
  \%.elf}).
\end{enumerate}
Il target \code{clean} rimuove tutti gli artefatti intermedi e le immagini
finali, per una ricompilazione pulita.

\textbf{Layout del flash} (il backing store di una U-proc): i blocchi
\code{0}--\code{30} contengono le pagine di \code{.text}/\code{.data}, il
blocco \code{31} la pagina di stack iniziale --- corrispondendo esattamente
alle 32 entry della Page Table privata.

\subsection{\code{phase3\_config\_machine.json} --- la configurazione dell'emulatore}
\begin{lstlisting}
{
    "boot": {
        "core-file": "build/MultiPandOS.core.uriscv",
        "load-core-file": true
    },
    "devices": {
        "flash0": { "enabled": true, "file": "testers/shell.uriscv" },
        "flash1": { "enabled": true, "file": "testers/echo.uriscv" },
        "flash2": { "enabled": true, "file": "testers/fibEleven.uriscv" },
        "flash3": { "enabled": true, "file": "testers/uname.uriscv" },
        "flash4": { "enabled": true, "file": "testers/sl.uriscv" },
        "flash5": { "enabled": true, "file": "testers/calc.uriscv" },
        "flash6": { "enabled": false, "file": "" },
        "flash7": { "enabled": false, "file": "" },
        "terminal0": { "enabled": true, "file": "term0.uriscv" }
    },
    "num-processors": 1,
    "num-ram-frames": 256,
    "symbol-table": { "asid": 64, "file": "build/MultiPandOS.stab.uriscv" },
    "tlb-floor-address": "0x80000000",
    "tlb-size": 16
}
\end{lstlisting}
Rispetto a \code{config\_machine.json} usato in Fase 1/2 (documento di Fase
1), tre differenze concrete che si ricollegano direttamente a ciò che
questo documento ha spiegato:
\begin{itemize}
\item \textbf{Sei device flash abilitati} (\code{flash0}--\code{flash5}),
  ciascuno puntato a un'immagine \code{testers/*.uriscv} generata dal
  Makefile appena visto. La \textbf{mappatura è posizionale}: il device
  \code{flashN} corrisponde all'ASID \code{N+1}, esattamente come calcolato
  da \code{flashOperation} (\code{devNo = asid - 1}) in
  \code{vmSupport.c}. \code{flash6}/\code{flash7} restano disabilitati:
  \code{UPROCMAX = 8} ne permetterebbe fino a 8, ma solo 6 programmi sono
  forniti.
\item \code{num-ram-frames: 256}, molto più dei 64 di Fase 1/2: deve
  ospitare i 32 frame riservati al Nucleus \emph{più} i 16 frame dello
  Swap Pool \emph{più} margine per gli stack degli 8 possibili
  \code{support\_t} (ciascuna con due stack da 500 word) --- la memoria
  virtuale ha bisogno di ben più fisico spazio di manovra del semplice
  Nucleus.
\item \code{tlb-floor-address: "0x80000000"}: a differenza di Fase 1/2
  (dove valeva \code{"0x0"}, nessuno spazio utente da proteggere), qui
  coincide esattamente con \code{KUSEG}: dice all'hardware TLB dove inizia
  lo spazio di indirizzamento \emph{utente} che deve essere tradotto
  tramite le Page Table private delle U-proc (indirizzi sotto questa soglia
  restano nello spazio kernel, mappato in modo diretto/identità).
\end{itemize}

\textbf{Da tenere bene a mente}: il flash \emph{non} è un filesystem
generico, ma il backing store grezzo di \emph{una singola} U-proc. Il Pager
lo legge a domanda (al verificarsi di un page fault) e vi riscrive le
pagine \emph{dirty} sfrattate; non c'è alcuna nozione di file, directory o
più programmi condivisi sullo stesso device.

% ============================================================
\section{Mappa riassuntiva della memoria (Fase 3)}
% ============================================================
\begin{center}
\begin{tabular}{@{}lll@{}}
\toprule
\textbf{Indirizzo} & \textbf{Costante} & \textbf{Contenuto} \\
\midrule
\code{0x20000000} & \code{RAMSTART} & inizio RAM; primi 32 frame riservati al Nucleus \\
\code{0x20020000} & \code{SWAP\_POOL\_START} & inizio dei 16 frame dello Swap Pool \\
\code{0x80000000} & \code{KUSEG} / \code{TEXT\_VADDR} & inizio spazio logico
  di una U-proc (\code{.text}) \\
\code{0x800000B0} & \code{UPROCSTARTADDR} & simbolo \code{\_\_start}, punto
  di ingresso reale (dopo l'header \code{.aout}) \\
\code{0xBFFFF000} & VPN \code{0xBFFFF} & pagina di stack logica (indice 31
  della Page Table) \\
\code{0xC0000000} & \code{USERSTACKTOP} & cima dello stack utente (valore
  iniziale di \code{sp}) \\
\bottomrule
\end{tabular}
\end{center}


\end{document}
